% !TeX root = ../../Book4.tex
% 纲要标题：### 11.2 截短消解与同调维数估计
% 纲要标签：b4:sec:1002
% 纲要位置：计划纲要.md，第 3471 行。

\section{截短消解与同调维数估计}
\label{b4:sec:1002}

高次导出群是否消失，可以转化为消解末端的合冲模是否投射。这个判据既判断给定消解能否截短，也说明局部环上的投射维数可以怎样检测。

% 纲要内容：
%
% * 截短消解及合冲模
% * 短正合列中的同调维数不等式
% * 第三卷局部投射维数所需的最低平移版本已在第9.5节证明；本节作一般化整理
% * 局部投射维数的讨论作为交换代数应用，可在学习第三卷后回读；不参与前两小节的证明
%

\subsection{截短消解及合冲模}
\label{b4:subsec:1002:01}

维数的定义只要求存在一条足够短的消解；下面的判据则作用于任意一条已经给定的消解。若 \(\operatorname{pd}_RM\le d\)，无论前 \(d\) 步选用了哪些投射满射，得到的第 \(d\) 个合冲模都投射，因而都能在此截短。内射方向有同样的结论；对于平坦维数，当前给定的各项只需平坦，不必投射。

\begin{proposition}\label{b4:prop:truncation-dimension}
设 \(R\) 为含幺环，\(M\) 为幺性左 \(R\) 模，\(d\ge1\) 为整数。给定幺性左 \(R\) 模正合列
\[
 0\longrightarrow K_d\longrightarrow P_{d-1}\longrightarrow\cdots
 \longrightarrow P_0\longrightarrow M\longrightarrow0,
\]
其中各 \(P_i\) 投射。则 \(\operatorname{pd}_RM\le d\) 当且仅当 \(K_d\) 投射。若幺性左 \(R\) 模序列 \(0\to M\to I^0\to\cdots\to I^{d-1}\to L^d\to0\) 正合且各 \(I^i\) 内射，则 \(\operatorname{id}_RM\le d\) 当且仅当 \(L^d\) 内射。若第一列的 \(P_i\) 仅要求平坦，则 \(\operatorname{fd}_RM\le d\) 当且仅当 \(K_d\) 平坦。
\end{proposition}
\begin{proof}
若末端对象具有所需性质，将它作为第 \(d\) 个投射项、内射项或平坦项，给定列本身就是长度至多 \(d\) 的相应消解。

反向令 \(K_0=M\)，把第一列拆成 \(0\to K_{j+1}\to P_j\to K_j\to0\)（\(0\le j<d\)）；这里 \(P_j\to K_j\) 是原微分到其像的满射，最后的核就是给定的 \(K_d\)。若各 \(P_j\) 投射，对任意幺性左模 \(N\) 依次用\cref{b4:thm:dimension-shift-basic}，得到
\[
 \operatorname{Ext}_R^1(K_d,N)\simeq\operatorname{Ext}_R^{d+1}(M,N).
\]
\(\operatorname{pd}_RM\le d\) 使右侧对所有 \(N\) 消失；\cref{b4:thm:ext-projective-injective}的一次判据便使这个给定的 \(K_d\) 投射。

内射列则令 \(L^0=M\)，依次取余核，拆成 \(0\to L^j\to I^j\to L^{j+1}\to0\)（\(0\le j<d\)）。第二变量的平移\cref{b4:prop:injective-dimension-shift}给出
\[
 \operatorname{Ext}_R^1(N,L^d)\simeq\operatorname{Ext}_R^{d+1}(N,M).
\]
若 \(\operatorname{id}_RM\le d\)，右侧对所有幺性左模 \(N\) 消失，一次 Ext 判据使 \(L^d\) 内射。

平坦版本仍拆第一列，但不使用 \(P_j\) 的投射性。对任意幺性右模 \(T\)，\cref{b4:thm:tor-flat-criterion}使 \(\operatorname{Tor}_i^R(T,P_j)=0\) 对每个 \(i>0\) 成立；\cref{b4:thm:tor-long-exact}于是逐步给出
\[
 \operatorname{Tor}_1^R(T,K_d)\simeq\operatorname{Tor}_{d+1}^R(T,M).
\]
若 \(\operatorname{fd}_RM\le d\)，右侧由\cref{b4:thm:dimension-vanishing}对所有 \(T\) 消失，同一平坦判据使 \(K_d\) 平坦。
\end{proof}

因此某个 \(d\) 步合冲模投射时，任何投射消解的 \(d\) 步合冲模都投射。这些合冲模不必同构，但加入投射直和项后同构，即稳定同构。具体地，两个投射消解的第一步给出短正合列 \(0\to K_1\to P_0\to M\to0\) 与 \(0\to K'_1\to P'_0\to M\to0\)，由\cref{b4:lem:schanuel}有
\[
 K_1\oplus P'_0\cong K'_1\oplus P_0.
\]
对第二步，记原消解给出的满射为 \(\pi_1:P_1\to K_1\)、\(\pi'_1:P'_1\to K'_1\)。分别补入恒等映射，得到
\[
\begin{aligned}
0&\longrightarrow K_2\longrightarrow P_1\oplus P'_0
 \xrightarrow{\pi_1\oplus1}K_1\oplus P'_0\longrightarrow0,\\
0&\longrightarrow K'_2\longrightarrow P'_1\oplus P_0
 \xrightarrow{\pi'_1\oplus1}K'_1\oplus P_0\longrightarrow0.
\end{aligned}
\]
恒等分量的核为零，所以两条新列的核仍分别是 \(K_2,K'_2\)；中间项都是投射模。利用第一步的同构将两个右端识别，再用 Schanuel 引理，得到
\[
 K_2\oplus P'_1\oplus P_0\cong K'_2\oplus P_1\oplus P'_0.
\]
同样，若在第 \(j\) 步已有 \(K_j\oplus Q\cong K'_j\oplus Q'\)，其中 \(Q,Q'\) 投射，就比较 \(P_j\oplus Q\to K_j\oplus Q\) 与 \(P'_j\oplus Q'\to K'_j\oplus Q'\)，恒等分量仍不改变核。再次应用该引理，便得到第 \(j+1\) 步的稳定同构。归纳可得投射模 \(Q,Q'\)，使 \(K_d\oplus Q\cong K'_d\oplus Q'\)。同调维数检测的是这些合冲模何时投射，并不要求它们作为模相互同构。

\subsection{短正合列中的同调维数不等式}
\label{b4:subsec:1002:02}

\begin{proposition}\label{b4:prop:dimension-short-exact}
设 \(R\) 为含幺环，\(0\to A\to B\to C\to0\) 为幺性左 \(R\) 模短正合列。以下所有维数均在 \(R\) 上计算。投射维数满足
\[
\begin{aligned}
 \operatorname{pd}B&\le\max\{\operatorname{pd}A,\operatorname{pd}C\},\\
 \operatorname{pd}A&\le\max\{\operatorname{pd}B,\operatorname{pd}C-1\},\\
 \operatorname{pd}C&\le\max\{\operatorname{pd}B,\operatorname{pd}A+1\}.
\end{aligned}
\]
平坦维数满足相同的三式。内射维数满足
\[
\begin{aligned}
 \operatorname{id}B&\le\max\{\operatorname{id}A,\operatorname{id}C\},\\
 \operatorname{id}A&\le\max\{\operatorname{id}B,\operatorname{id}C+1\},\\
 \operatorname{id}C&\le\max\{\operatorname{id}B,\operatorname{id}A-1\}.
\end{aligned}
\]
这里维数取值及 \(\pm\infty\) 的运算遵循\cref{b4:def:homological-dimensions}中的约定。
\end{proposition}
\begin{proof}
若某一右端为 \(\infty\)，该不等式直接成立。右端为负的情形由短列的正合性迫使被估计的模为零，结论也直接成立。以下考虑非负有限上界；零次 Ext、Tor 分别按 Hom、张量积解释。

固定幺性左模 \(N\)，\cref{b4:thm:ext-long-exact}在第一变量给出五项片段
\[
\begin{aligned}
 \operatorname{Ext}^{i-1}(A,N)&\longrightarrow\operatorname{Ext}^i(C,N)
 \longrightarrow\operatorname{Ext}^i(B,N)\\
 &\longrightarrow\operatorname{Ext}^i(A,N)
 \longrightarrow\operatorname{Ext}^{i+1}(C,N)
 \qquad(i\ge1).
\end{aligned}
\]
估计 \(B\) 时，取 \(i>\max\{\operatorname{pd}A,\operatorname{pd}C\}\)，它两侧的同次数 Ext 均为零，故 \(\operatorname{Ext}^i(B,N)=0\)。估计 \(A\) 时，条件
\(i>\max\{\operatorname{pd}B,\operatorname{pd}C-1\}\)
同时给出 \(i>\operatorname{pd}B\) 与 \(i+1>\operatorname{pd}C\)，使夹住 \(\operatorname{Ext}^i(A,N)\) 的两项消失。估计 \(C\) 时，条件
\(i>\max\{\operatorname{pd}B,\operatorname{pd}A+1\}\)
则给出 \(i>\operatorname{pd}B\) 与 \(i-1>\operatorname{pd}A\)，使 \(\operatorname{Ext}^{i-1}(A,N)\) 与 \(\operatorname{Ext}^i(B,N)\) 消失。这说明三个最大值中的次数偏移恰好来自 \(i+1\) 和 \(i-1\)。以上对每个 \(N\) 成立，\cref{b4:thm:dimension-vanishing}便给出三条投射维数不等式。

平坦情形固定幺性右模 \(T\)，\cref{b4:thm:tor-long-exact}给出
\[
\begin{aligned}
 \operatorname{Tor}_{i+1}(T,C)&\longrightarrow\operatorname{Tor}_i(T,A)
 \longrightarrow\operatorname{Tor}_i(T,B)\\
 &\longrightarrow\operatorname{Tor}_i(T,C)
 \longrightarrow\operatorname{Tor}_{i-1}(T,A)
 \qquad(i\ge1).
\end{aligned}
\]
\(B\) 仍被 \(A,C\) 的同次数项夹住；\(A\) 被 \(C\) 的 \(i+1\) 次与 \(B\) 的 \(i\) 次项夹住；\(C\) 被 \(B\) 的 \(i\) 次与 \(A\) 的 \(i-1\) 次项夹住。消失所需的次数条件与投射情形相同，故得到同样的三式。

内射情形固定幺性左模 \(N\)，改用第二变量的片段
\[
\begin{aligned}
 \operatorname{Ext}^{i-1}(N,C)&\longrightarrow\operatorname{Ext}^i(N,A)
 \longrightarrow\operatorname{Ext}^i(N,B)\\
 &\longrightarrow\operatorname{Ext}^i(N,C)
 \longrightarrow\operatorname{Ext}^{i+1}(N,A)
 \qquad(i\ge1).
\end{aligned}
\]
估计 \(B\) 仍只用 \(A,C\) 的同次数项。估计 \(A\) 却要令 \(i>\operatorname{id}B\)、\(i-1>\operatorname{id}C\)，所以其上界是 \(\max\{\operatorname{id}B,\operatorname{id}C+1\}\)；这时左侧的 \(\operatorname{Ext}^{i-1}(N,C)\) 与右侧的 \(\operatorname{Ext}^i(N,B)\) 都为零，正合性使 \(\operatorname{Ext}^i(N,A)=0\)。估计 \(C\) 则要令 \(i>\operatorname{id}B\)、\(i+1>\operatorname{id}A\)，上界便是 \(\max\{\operatorname{id}B,\operatorname{id}A-1\}\)。因此两处偏移与投射情形不同；对所有 \(N\) 使用内射消失判据即得结论。
\end{proof}

这些不等式也给出精确的一步下降。在 \(0\to K\to P\to M\to0\) 中，若 \(P\) 投射且 \(\operatorname{pd}M=d\ge1\) 有限，则 \(P\ne0\)，故中间项的维数固定为 \(\operatorname{pd}P=0\)。两个不同的不等式于是能分别给出 \(K\) 的上界与下界。第二个投射维数不等式给出
\[
 \operatorname{pd}K\le\max\{\operatorname{pd}P,\operatorname{pd}M-1\}
 =\max\{0,d-1\}=d-1.
\]
第三个投射维数不等式则给出
\[
 d=\operatorname{pd}M\le\max\{\operatorname{pd}P,\operatorname{pd}K+1\}
 =\max\{0,\operatorname{pd}K+1\}.
\]
因 \(d\ge1\)，后一式迫使 \(\operatorname{pd}K\ge d-1\)，从而 \(\operatorname{pd}K=d-1\)。特别地，\(d=1\) 时两式分别给出 \(\operatorname{pd}K\le0\) 与 \(\operatorname{pd}K\ge0\)，所以 \(K\) 是非零投射模。若 \(d=0\)，短列分裂，\(K\) 投射或为零；不能强行套用数值 \(d-1=-1\)。

\subsection{交换代数应用：局部投射维数}
\label{b4:subsec:1002:03}

一般环上的消失判据需要对所有测试模作检验。能否只选一个剩余域作测试，取决于消解在这个域上能保留什么信息：若微分仍非零，单个同调群为零只表示该处的核等于像，并不能说明对应自由项为零。

\ProseParagraphBreak
% 跨卷引用目标：b3:prop:minimal-resolution-existence（17.1.4）。
% 跨卷引用目标：b3:prop:minimal-resolution-betti（17.1.5）。
设 \((R,\mathfrak m,k)\) 为交换 Noether 局部环，\(M\ne0\) 为有限生成的 \(R\) 模。取各项均为有限秩的极小自由消解 \(F_\bullet\to M\)，即每个正次微分满足 \(\mathrm{d}_i(F_i)\subseteq\mathfrak mF_{i-1}\)。因为 \(k=R/\mathfrak m\)，微分矩阵与 \(k\) 张量后全部变为零；向 \(k\) 作 Hom 时，\(\mathfrak m\) 的作用也为零。这样的消解由第三卷命题17.1.4构造，命题17.1.5据此给出
\[
 \operatorname{Tor}_i^R(k,M)\simeq k\otimes_R F_i,\qquad
 \operatorname{Ext}_R^i(M,k)\simeq\operatorname{Hom}_R(F_i,k).
\]
这里 \(R\) 交换，\(k\) 同时视为右 \(R\) 模，Tor 的平衡计算使我们可以用第二变量的自由消解 \(F_\bullet\)。两个复形已无非零微分，故每个整数 \(i\ge0\) 的同调就是该项本身，模 \(M\) 的第 \(i\) 个 Betti 数\footnote{贝蒂（Enrico Betti，1823-1892），意大利数学家，研究代数与拓扑。}为
\[
 \beta_i^R(M)=\operatorname{rank}_R F_i
 =\dim_k\operatorname{Tor}_i^R(k,M).
\]
% 跨卷引用目标：b3:thm:nakayama（3.2.3）。
每个 \(F_i\) 有限生成，第三卷定理3.2.3的 Nakayama 引理使 \(k\otimes_RF_i=F_i/\mathfrak mF_i=0\) 等价于 \(F_i=0\)。还须检查零项之后会不会重新出现非零项。若 \(F_i=0\)，则 \(\mathrm{d}_{i+1}:F_{i+1}\to F_i\) 的目标为零，因而整个 \(F_{i+1}\) 都是其核；消解在 \(F_{i+1}\) 处正合，把这个核识别为 \(\mathrm{d}_{i+2}\) 的像；极小性又使该像包含在 \(\mathfrak mF_{i+1}\) 中。因此
\[
F_{i+1}=\ker\mathrm{d}_{i+1}=\operatorname{im}\mathrm{d}_{i+2}
\subseteq\mathfrak mF_{i+1}.
\]
由 \(F_{i+1}\) 有限生成，再用 Nakayama 引理得到 \(F_{i+1}=0\)。递推后所有更高项都为零；一般正合复形在一个零项后仍可接上非零的两项恒等复形，极小性排除了这种情形。

\ProseParagraphBreak
% 跨卷引用目标：b3:thm:local-pd-vanishing（17.2.3）。
对整数 \(d\ge0\)，若 \(\operatorname{Tor}_{d+1}^R(k,M)=0\)，便有 \(F_{d+1}=0\)，后续各项也为零，因而这条极小消解的长度不超过 \(d\)。反过来，长度不超过 \(d\) 的投射消解也是平坦消解，所以 \(\operatorname{pd}_RM\le d\) 使所有高于 \(d\) 次的 Tor 消失，尤其使这个剩余域测试消失。Ext 的计算同样检测 \(F_i\) 是否为零。因此
\[
\begin{aligned}
\operatorname{pd}_RM
&=\sup\{\,i\ge0:\operatorname{Tor}_i^R(k,M)\ne0\,\}\\
&=\sup\{\,i\ge0:\operatorname{Ext}_R^i(M,k)\ne0\,\},
\end{aligned}
\]
这也解释了第三卷定理17.2.3中只检验剩余域的判据。这里的局部性、极小性及各项有限生成性共同保证检测有效。

\ProseParagraphBreak
例如，设 \(k\) 为域。在 \(R=k[x,y]_{(x,y)}\) 上，剩余域的 Koszul 消解为
\[
 0\longrightarrow R\xrightarrow{\binom{-y}{x}}R^2
 \xrightarrow{(x\ \ y)}R\longrightarrow k\longrightarrow0.
\]
它极小，故 \(\operatorname{Tor}_2^R(k,k)=k\ne0\)，同时长度二给出高次消失，从而 \(\operatorname{pd}_Rk=2\)。相反，\cref{b4:exm:infinite-dimension-dual-numbers}的每一步都保留一个自由项，有限生成并不意味着有限投射维数。在有限投射维数情形，第三卷第17.3节的 Auslander--Buchsbaum 公式\footnote{奥斯兰德（Maurice Auslander，1926-1994），美国数学家，研究同调代数与表示论；布赫斯鲍姆（David Alvin Buchsbaum，1929-2021），美国数学家，研究同调代数及交换代数中的自由消解。}进一步将投射维数与模和底环的深度联系起来。
